\section*{Hazard Analysis Appendix}

The following hazards were identified and mitigated by design and/or
process, covering personnel safety during integration, launch and
recovery operations. The complete hazard analysis is presented in
Table~\ref{tab:hazards}:

\vspace{0.3cm}

\begin{table}[htbp]
\centering
\small
\setlength{\tabcolsep}{6pt}
% \tabularx exige a largura total (\textwidth) e usa a coluna do tipo 'X' para autoajuste
\begin{tabularx}{\textwidth}{p{3.5cm} X X} 
\toprule
\textbf{Hazard ID \& Title} & \textbf{Design Mitigation} & \textbf{Process Mitigation} \\
\midrule

\textbf{Hazard 1} \newline Battery thermal runaway &
BMS with multiple protections covering overcharge, overdischarge, short-circuit, and overcurrent conditions. &
Cells cycle-tested before integration. Battery charged only with approved charger, monitored during pre-launch, and RBF pin cuts all energy during transport. \\
\addlinespace

\textbf{Hazard 2} \newline Stored energy during integration &
Two independent mechanical kill switches (SW1/SW2) wired in parallel (NC) on main power line. RBF pin (SW3) in series cuts all energy when inserted. &
Pre-integration verification procedure checks kill switch activation. RBF removed only after installation in the deployer. \\
\addlinespace

\textbf{Hazard 3} \newline Sharp edges of SRAB wings &
Final geometry features rounded leading edges. &
Team wears handling gloves during integration. Recovery team briefed on wing handling. \\
\addlinespace

\textbf{Hazard 4} \newline RF interference with other teams &
Dedicated packet protocol with team ID, a reserved 915~MHz channel, and local storage for full data backup. &
Coordinated frequency allocation during operations. \\
\addlinespace

\textbf{Hazard 5} \newline Uncontrolled descent &
Passive deployment with no actuators, two-wing configuration for predictable recovery, keeping nominal impact energy low (17.8~J). &
Validated through three iterative simulation-and-test campaigns. \\

\bottomrule
\end{tabularx}
\caption{Hazard analysis and mitigations.}
\label{tab:hazards}
\end{table}
